%----------------------------------------------------------------------------
\chapter{\bevezetes}
%----------------------------------------------------------------------------

A sakk évtizedek óta a mesterséges intelligencia kutatásának egyik kedvelt terepe: szabályai pontosan definiáltak, állapottere azonban olyan hatalmas, hogy a teljes játékfa bejárása gyakorlatilag lehetetlen. Claude Shannon már 1950-ben megfogalmazta a számítógépes sakk alapproblémáját, és megbecsülte, hogy egy átlagos játszma lehetséges folytatásainak száma nagyságrendileg $10^{120}$~\cite{Shannon50}. Egy erős sakkmotor ezért nem brute-force kereséssel, hanem jól megválasztott vágási és értékelési heurisztikákkal dolgozik.

A szakdolgozatom célja egy saját fejlesztésű sakkmotor megtervezése és megvalósítása volt, amely
\begin{itemize}
	\item a klasszikus (FIDE) szabályok szerint játszik,
	\item az iparági szabványnak tekinthető UCI (Universal Chess Interface) protokollon keresztül más programokkal is együttműködik,
	\item saját, interaktív grafikus felhasználói felülettel rendelkezik,
	\item és úgy épül fel, hogy új, nem szabványos (tündérsakk) bábukkal is bővíthető legyen.
\end{itemize}

A tündérsakk-kiterjesztés bemutatására az Amazon bábut választottam, amely a vezér és a huszár lépésmódját egyesíti. Egy ilyen bábu bevezetése nemcsak a lépésgenerátort érinti, hanem a cserék kiértékelését, a lépésrendezést és az állásértékelést is, így jól megmutatja, mennyire általános egy motor architektúrája.

A munka két forráskód-tárolóban készült. A backend a motor konzolos változata, ebben fut a UCI és a robot-robot elleni tömeges mérés. A GUI projekt ugyanezt a motort egy grafikus felülettel egészíti ki. A fejlesztés 2025 decemberében indult, a motor 2026 áprilisától kezdve mért, verziózott lépésekben fejlődött: minden módosított keresőt egy korábbi verzió ellen több ezer játszmában teszteltem, és csak a bizonyítottan erősebb változatot tartottam meg.

\section*{A dolgozat felépítése}

\Aref{chap:rules}. fejezet bemutatja a vonatkozó FIDE-szabályokat és a tündérsakk bábuk alapelveit. \Aref{chap:background}. fejezet a sakkmotorok szakirodalmi hátterét foglalja össze. \Aref{chap:architecture}. fejezet a szoftverarchitektúrát és a UCI kommunikációs réteget, \aref{chap:movegen}. fejezet a táblareprezentációt és a lépésgenerálást ismerteti. \Aref{chap:search}. fejezet a keresőmotor összes funkcióját és optimalizációját tárgyalja, verziók szerint. \Aref{chap:gui}. fejezet a grafikus felületet, \aref{chap:amazon}. fejezet az Amazon bábuval bővített játékváltozatot mutatja be. \Aref{chap:measurements}. fejezetben mérésekkel igazolom a lépésgenerálás helyességét és hatékonyságát, valamint a motor fejlődését. A dolgozatot az eredmények értékelése és a továbbfejlesztési lehetőségek zárják.
